\begin{figure}[htbp]
\centering
\begin{tikzpicture}[>=Latex,font=\small,
 state/.style={draw,rounded corners,fill=blue!4,minimum height=10mm,
               minimum width=21mm,inner sep=4pt}]
 \node[state] (u) at (0,1.65) {$u_h$};
 \node[state] (first) at (3.7,1.65) {$\partial_{x_i}u_h$};
 \node[state] (second) at (8,1.65) {$\partial_{x_j}\partial_{x_i}u_h$};
 \draw[->,thick] (u)--node[above] {$\partial_{x_i}$}(first);
 \draw[->,black!55,dashed,thick] (first)--node[above] {$\partial_{x_j}$}(second);
 \node[state,fill=teal!6] (projected) at (3.7,-.25) {$\Pi\partial_{x_i}u_h$};
 \node[state,fill=teal!6] (recursive) at (8,-.25)
    {$\Pi\partial_{x_j}(\Pi\partial_{x_i}u_h)$};
 \draw[->,teal!80!black,thick] (first)--node[left] {project: $\Pi$}(projected);
 \draw[->,teal!80!black,thick] (projected)--node[above]
    {$\Pi\partial_{x_j}$}(recursive);
 \node[align=left,anchor=east] at (1.9,-.25)
    {implemented route\\back into $V^k(K)$};
 \node[align=center] at (8,-1.55)
    {generally $\neq\Pi\partial_{x_j}\partial_{x_i}u_h$};
\end{tikzpicture}
\caption{Exact differentiation and the implemented recursive construction on a
curved element, with $\Pi=\Pi_K^k$. The first physical derivative generally
leaves $V^k(K)$. Projecting it back changes the input to the next derivative.
The lower route therefore differs from differentiating twice and projecting
only at the end. Each lower state is represented in the same finite element
space; no second derivative of the geometry map is needed by this route.}
\label{fig:projected-derivative-route}
\end{figure}
